% Generated draft: requires mathematical and source review.
\documentclass[11pt]{article}
\usepackage[margin=25mm]{geometry}
\usepackage{fontspec}
\setmainfont{Latin Modern Roman}
\usepackage{amsmath,amssymb,mathtools,unicode-math}
\setmathfont{Latin Modern Math}
\usepackage{enumitem,needspace,xurl,hyperref}
\setlength{\parindent}{0pt}
\setlength{\parskip}{5pt}
\setlength{\emergencystretch}{3em}
\newcommand{\N}{\mathbb{N}}
\newcommand{\Z}{\mathbb{Z}}
\newcommand{\Q}{\mathbb{Q}}
\newcommand{\R}{\mathbb{R}}
\newcommand{\C}{\mathbb{C}}
\newcommand{\F}{\mathbb{F}}
\newcommand{\A}{\mathbb{A}}
\newcommand{\PP}{\mathbb{P}}
\newcommand{\E}{\mathbb{E}}
\newcommand{\eps}{\varepsilon}
\newcommand{\dd}{\,\mathrm{d}}
\begin{document}
\section*{UnsolvedMath ID: 30000184}
\textbf{Source problem code:} catalogue entry
\subsection{Definitions and mathematical statement}
\textbf{Retained source question.} Is this family \((B;A_\lambda)\) of pairwise nonisomorphic subgroup embeddings the first family which occurs? A subgroup embedding can be represented by a pair \((B;A)\) with \(A\le B\). Two such pairs are isomorphic if an isomorphism of ambient groups \(\phi:B\to B'\) satisfies \(\phi(A)=A'\). For fixed \(B\), pairwise nonisomorphism means the \(A_\lambda\) belong to distinct \(\operatorname{Aut}(B)\)-orbits. This defines the ordinary embedding terminology but does not reconstruct the missing family.\par The ambient group or module \(B\), subgroups \(A_\lambda\), parameter set, prime and exponent, and the ordering measured by \emph{first} are absent. The supplied assessment mentions the classical \(\Z/p^n\) setting, but does not say whether minimality concerns exponent, ambient type, ambient order, subgroup type, or another complexity. Nor is \emph{family} defined as a parameterized construction across primes or rings, a positive-dimensional family over an infinite field, or a finite collection for fixed data. These possibilities cannot be identified silently; for a fixed finite group the number of subgroups is finite. Thus no precise minimality proposition can be extracted from the surviving question alone.\par \textbf{Conflicting supplied status notes.} The triage checked 2026-08-21 labels the item open and says checked literature did not settle minimality, despite finite-type lower bounds and related examples. The later statement-verification field instead reports that the official report's final theorem answers this very minimal-family question positively, labels the statement \texttt{unrecoverable}, and expressly declines to promote a separate clean open statement. The theorem and the family are not reproduced. This draft retains that reported positive answer and the conflicting dated triage as provenance; it does not invent an open conjecture, a theorem statement, or a recovered family.
\subsection{Short English statement}
The question concerns minimality of a particular family of nonisomorphic subgroup embeddings, but the family and the meaning of minimality are missing. A later verification note reports a positive answer in the source itself, conflicting with the dated open-status triage. No separate clean open problem is reconstructed.
\subsection{Sources}
\begin{itemize}\item Local UnsolvedMath registry, permanent ID 30000184, code catalogue entry; \emph{Minimal Families of Nonisomorphic Subgroup Embeddings}, Oberwolfach Reports Open Problems. Statement status \texttt{unrecoverable}; no clean statement supplied.\item \emph{Representation Theory of Finite-Dimensional Algebras (2005)}, Oberwolfach Reports, DOI 10.4171/owr/2005/06, pp. 319--406. \url{https://doi.org/10.4171/owr/2005/06}. Proposed year 2005; no proposer named.\item Supplied literature label \emph{Schmidmeier, subgroup embeddings, original OWR report}: \url{https://doi.org/10.4171/owr/2005/06}. The final theorem reported to answer the question positively is not included in this record.\item Literature triage checked 2026-08-21, registry status \texttt{open}; record updated 2026-09-16. The conflicting later verification claim was not independently checked here.\end{itemize}
\end{document}
